% ── §1.6: Συνεισφορά στην Επιστήμη ──────────────────────────

\section{Συνεισφορά στην Επιστήμη}
\label{sec:contribution}

Η παρούσα εργασία συνεισφέρει στο επιστημονικό πεδίο σε τέσσερα συμπληρωματικά επίπεδα, που μαζί καλύπτουν ένα κενό που εντοπίστηκε στη βιβλιογραφία: η πλειοψηφία των υπαρχουσών εργασιών αντιμετωπίζει ένα από τα τρία προβλήματα (ιδιωτικότητα, κίνητρα, ή δικαιοσύνη) σε απομόνωση, χωρίς να τα ενσωματώνει σε ενιαίο λειτουργικό σύστημα αξιολογημένο σε ρεαλιστικό AAL dataset. Τα σημεία συνεισφοράς είναι:

\begin{enumerate}
    \item \textbf{Ολοκληρωμένο σύστημα FL για HAR σε AAL:} Σε αντίθεση
    με εργασίες που αντιμετωπίζουν μεμονωμένες πτυχές του προβλήματος,
    η παρούσα εργασία προτείνει ένα ενιαίο πλαίσιο που συνδυάζει
    ιδιωτικότητα, κίνητρα και δίκαιη αξιολόγηση.

    \item \textbf{Εφαρμογή παιγνίου Stackelberg σε AAL:} Η προσαρμογή
    του μηχανισμού Stackelberg στις ιδιαιτερότητες των AAL περιβαλλόντων
    (ετερογένεια χρηστών, ευαισθησία δεδομένων) αποτελεί πρωτότυπη
    συνεισφορά.

    \item \textbf{Συγκριτική αξιολόγηση μεθόδων συνεισφοράς:} Η
    συστηματική σύγκριση volume-based, LOO και gradient-based Shapley
    μεθόδων στο πλαίσιο HAR παρέχει νέα εμπειρικά δεδομένα για την
    επιλογή κατάλληλης μεθόδου ανάλογα με τους υπολογιστικούς πόρους.

    \item \textbf{Αξιολόγηση σε Non-IID AAL δεδομένα:} Η χρήση του
    AAL dataset \cite{AAL_UCI2016} παρέχει μια πιο ρεαλιστική
    αξιολόγηση σε σχέση με τα συνήθη benchmark datasets ελεγχόμενου
    περιβάλλοντος.
\end{enumerate}

Συνολικά, η εργασία αποτελεί ένα εμπειρικά θεμελιωμένο πλαίσιο για την ανάπτυξη FL συστημάτων HAR που σέβονται την ιδιωτικότητα, κινητοποιούν τη συμμετοχή και αξιολογούν δίκαια τη συνεισφορά — τρία αλληλένδετα κριτήρια που απαιτούν κοινή αντιμετώπιση για πρακτική εφαρμογή σε AAL περιβάλλοντα.
